\section{Solutions}\label{sec:13-next-steps:04-solutions:1}

\textbf{Socratic questions, looking back and forward.}

\begin{enumerate}
\def\labelenumi{\arabic{enumi}.}
\tightlist
\item
  \emph{Every group, ring, module, and path algebra in this book was built by hand, field by field. Having now seen Mathlib's \texttt{class}-based versions in §2, was the from-scratch work wasted effort, now that a library does it automatically?} No --- Mathlib automates exactly the parts that are safe to automate \emph{once the underlying data is already understood}. Nothing in §2 replaces knowing what a \texttt{Group} axiom actually demands; it only removes the bookkeeping of re-deriving that knowledge every time a new carrier type shows up.
\item
  \emph{Of the five projects sketched in §3, which one sounds least comfortable to attempt right now --- and is that discomfort a reason to avoid it, or a reason to pick it?} This book has no answer to give here; noticing which gap feels least settled is itself the most useful outcome of reaching this page, and that gap is usually exactly where a project pays off most.
\item
  \emph{This book verified every Lean snippet against a real toolchain rather than merely describing what should happen. Now that the training wheels of a guided worked example are gone, what is the equivalent habit to carry forward alone?} Trying the cheap thing first, reading why it failed, and checking the goal state after every step rather than only at the end --- the same loop Chapter 4 introduced, now without a book supplying the next line.
\end{enumerate}

Full worked solutions to every chapter's exercises are in the
